% ---
% id: statistics/binomial-box/thm:box-construction
% kind: thm
% title: "Properties of the constructed box"
% summary: >-
%   For any admissible endpoint rule (count bounds meeting binomial tail
%   budgets), the acceptance box has exact rational bounds and contains every
%   law within the honest envelope's bounds. It is non-empty and compact and,
%   by construction, an admissible acceptance set
%   ([[meat-core/def:admissible-acceptance-set]]).
% prominence: main
% status: implemented
% uses: [statistics/def:taxonomy, statistics/thm:geometry-validation, def:endpoint-rule, lem:guard, def:construction, statistics/def:envelope, statistics/def:conservative-evaluation]
% mentions: [meat-core/def:admissible-acceptance-set]
% assumes: [A-ALPHABET]
% refs: []
% provenance: newly-derived
% implements: certiqs_key.statistics.construction:construct
% lean:
% ---

Whatever admissible rule produced the endpoints, the output is a valid
acceptance set that accepts every honest law of the envelope: the guard, not
the rule, places the envelope inside the box.

\begin{theorem}\label{thm:box-construction}
On every admissible input (Definition~\ref{def:endpoint-rule}) and for every
admissible endpoint rule $R$, under a \tref{conservative evaluation}
(Definition~\ref{def:conservative-evaluation}):
\begin{enumerate}[label=(\roman*),leftmargin=2.2em]
\item $\mathsf{Acc}_R$ is well defined, and all of its output data are exact
rationals: each $q^{\pm}(\hat c)$ is either an integer multiple of $1/n$ or one
of $0,1$, and the entries of $B$ lie in $\{0,\pm1\}$;
\item $0\le\qlo\le\qhi\le\mathbf{1}$ componentwise, so the pair
$(\qlo,\qhi)$ meets the conditions of the canonical box form of
Definition~\ref{def:taxonomy};
\item $\mathsf{H}\subseteq\acSet$, where $\mathsf{H}$ is the
\tref{envelope set} \eqref{eq:envelope-set} of the input envelope: every
honest law the envelope admits is accepted, and for the degenerate envelope
this reads $\honest\in\acSet$. The envelope set is non-empty; consequently
$\acSet\neq\varnothing$ and, by Theorem~\ref{thm:geometry-validation}\,(iii),
$\sum_{\hat c}q^{-}(\hat c)\le1\le\sum_{\hat c}q^{+}(\hat c)$;
\item $\acSet$ is a non-empty, closed, convex, compact subset of
$\simplex{\Chatset}$; \xref{meat-core/def:admissible-acceptance-set} holds for the output of
$\mathsf{Acc}_R$ by construction, for every admissible input and every
admissible rule.
\end{enumerate}
\end{theorem}

\begin{proof}
\emph{(i).} Well-definedness: $R$ returns integers in $\{0,\dots,n\}$
(Definition~\ref{def:endpoint-rule}\,(a)), so the guarded endpoints
\eqref{eq:guard} exist and are unique, and
\eqref{eq:boxbounds}--\eqref{eq:accset} are explicit formulas in them. The
rationality statement is immediate: $k^{\pm}(\hat c)$ are integers in
$\{0,\dots,n\}$ and $n\ge1$.

\emph{(ii).} For $\hat c\in\Chatcon$, \eqref{eq:ordering} gives $0\le
k^{-}(\hat c)\le k^{+}(\hat c)\le n$, hence $0\le q^{-}(\hat c)\le q^{+}(\hat
c)\le1$. For $\hat c\notin\Chatcon$, $0\le0\le1\le1$.

\emph{(iii).} Let $\mathbf{p}\in\mathsf{H}$, let $\hat c\in\Chatcon$ and
write $p^{\mp}=p^{\mp}_{\hat c}$. By \eqref{eq:ordering} and
\eqref{eq:envelope-set},
$q^{-}(\hat c)=k^{-}(\hat c)/n\le\lfloor np^{-}\rfloor/n\le np^{-}/n=p^{-}
\le p(\hat c)$ and
$q^{+}(\hat c)=k^{+}(\hat c)/n\ge\lceil np^{+}\rceil/n\ge np^{+}/n=p^{+}
\ge p(\hat c)$, so $q^{-}(\hat c)\le p(\hat c)\le q^{+}(\hat c)$. For
$\hat c\notin\Chatcon$ the bounds are $0\le p(\hat c)\le1$, which hold because
$\mathbf{p}\in\simplex{\Chatset}$. Since also
$\mathbf{p}\in\simplex{\Chatset}$, all defining conditions of
\eqref{eq:accset} are met and $\mathbf{p}\in\acSet$; hence
$\mathsf{H}\subseteq\acSet$. For the degenerate envelope
$\mathsf{H}=\{\honest\}$ (Definition~\ref{def:envelope}).

Non-emptiness of $\mathsf{H}$: by Definition~\ref{def:envelope} the pair
$({\honest}^{-},{\honest}^{+})$ is a canonical box form with denotation
$\mathsf{H}$ whose bounds satisfy \eqref{eq:boxfeas}, so
Theorem~\ref{thm:geometry-validation}\,(iii), by which \eqref{eq:boxfeas} and
non-emptiness are equivalent, gives $\mathsf{H}\neq\varnothing$.
Hence $\acSet\supseteq\mathsf{H}$ is non-empty. The box sums follow from the
same equivalence applied to the non-empty $\acSet$, a canonical box by (ii).

\emph{(iv).} Theorem~\ref{thm:geometry-validation}\,(i) applied to the geometry of (ii),
whose \tref{validation predicate} holds by (iii).
\end{proof}
